\section{A tétel lényege és vizsgaszintű vázlata}

\begin{summarybox}
A tétel az algoritmusok elméleti határait vizsgálja. A fő kérdés nem az, hogy egy konkrét program gyors-e, hanem az, hogy egy pontosan megfogalmazott problémára létezik-e egyáltalán algoritmus, amely minden bemeneten helyes választ ad. Ehhez a problémákat nyelvekként kódoljuk, majd megkülönböztetjük az eldönthető, felismerhető és eldönthetetlen nyelveket. A tétel végén megjelenik a bonyolultságelméleti nézőpont is: az eldönthető problémák között különösen fontosak a polinomiális időben eldönthető problémák.
\end{summarybox}

Vizsgán a tétel jól felépíthető a következő sorrendben:

\begin{enumerate}
  \item az algoritmikus eldönthetőség motivációja: igen/nem kérdések és nyelvek;
  \item Church-tézis: az intuitív algoritmusfogalom és a formális modellek kapcsolata;
  \item rekurzív és parciálisan rekurzív függvények;
  \item rekurzív, illetve rekurzívan felsorolható nyelvek;
  \item felismerés és eldöntés különbsége;
  \item az $\mathrm{R}$, $\mathrm{RE}$, $\mathrm{coR}$ és $\mathrm{coRE}$ osztályok kapcsolata;
  \item nevezetes eldönthetetlen problémák, különösen a megállási probléma;
  \item redukciós gondolat: eldönthetetlenséget gyakran ismert eldönthetetlen problémából vezetünk le;
  \item időbonyolultság és polinomiális idejű algoritmusok;
  \item a tétel végén az összkép: nem minden pontos probléma algoritmizálható, és nem minden algoritmizálható probléma hatékonyan oldható meg.
\end{enumerate}

A legfontosabb vizsgamondat: egy nyelv akkor rekurzív, ha létezik hozzá minden bemeneten megálló döntő algoritmus; akkor rekurzívan felsorolható, ha az igen esetek felismerhetők véges idő alatt, de a nem esetekben a gép akár örökké is futhat. Az $\mathrm{R}$ osztály valódi részhalmaza az $\mathrm{RE}$ osztálynak, a megállási probléma pedig tipikus példa olyan problémára, amely felismerhető, de nem eldönthető.

\begin{center}
\begin{tabularx}{\textwidth}{L{22mm}L{34mm}L{38mm}X}
\toprule
\textbf{Osztály} & \textbf{Jelentés} & \textbf{Gépi jellemzés} & \textbf{Megjegyzés} \\
\midrule
$\mathrm{R}$ & Rekurzív, eldönthető nyelvek & Van olyan Turing-gép, amely minden bemeneten megáll, és helyes igen/nem választ ad. & Ezekre teljes algoritmus létezik. Ha $L\in\mathrm{R}$, akkor $\overline L\in\mathrm{R}$ is. \\
$\mathrm{RE}$ & Rekurzívan felsorolható, felismerhető nyelvek & Van olyan Turing-gép, amely $w\in L$ esetén elfogadással megáll; $w\notin L$ esetén elutasíthat vagy futhat örökké. & Félig eldönthető osztály. Az igen esetek tanúsíthatók véges futással. \\
$\mathrm{coR}$ & Az $\mathrm{R}$ komplementerosztálya & $L\in\mathrm{coR}$, ha $\overline L\in\mathrm{R}$. & Mivel $\mathrm{R}$ zárt komplementerre, $\mathrm{coR}=\mathrm{R}$. \\
$\mathrm{coRE}$ & Az $\mathrm{RE}$ komplementerosztálya & $L\in\mathrm{coRE}$, ha $\overline L\in\mathrm{RE}$. & A nem esetek félig eldönthetők. Általában $\mathrm{RE}\neq\mathrm{coRE}$. \\
\bottomrule
\end{tabularx}

\end{center}

\section{Algoritmikus eldönthetőség}

\subsection{Problémák mint igen/nem kérdések}

Az algoritmikus eldönthetőség olyan problémákkal foglalkozik, amelyekre a válasz igen vagy nem. Például:

\begin{itemize}
  \item prímszám-e egy adott természetes szám;
  \item elfogad-e egy adott véges automata egy adott szót;
  \item megáll-e egy adott program egy adott bemeneten;
  \item üres-e egy adott Turing-gép által felismert nyelv.
\end{itemize}

Az ilyen problémákat nyelvekként szoktuk formalizálni. Rögzítünk egy véges ábécét $\Sigma$, és minden objektumot véges szóként kódolunk. Egy probléma ekkor egy $L\subseteq\Sigma^*$ nyelvvel adható meg:
\[
  w\in L \quad\Longleftrightarrow\quad
  \text{a } w \text{ kódolású példányra az igen válasz helyes.}
\]

Például a megállási probléma nyelve:
\[
  \mathrm{HALT}=\{\langle M,w\rangle\mid
  M \text{ Turing-gép megáll a } w \text{ bemeneten}\}.
\]
Itt $\langle M,w\rangle$ azt jelenti, hogy a $M$ gép leírását és a $w$ bemenetet egyetlen véges szóként kódoljuk.

\subsection{Eldönthető probléma}

Egy $L\subseteq\Sigma^*$ nyelv \textbf{eldönthető}, ha létezik olyan algoritmus, amely minden $w\in\Sigma^*$ bemenetre véges idő után megáll, és helyesen megmondja, hogy $w\in L$ vagy $w\notin L$.

Turing-géppel megfogalmazva: $L$ eldönthető, ha van olyan Turing-gép, amely minden bemeneten megáll, és pontosan az $L$-beli szavakat fogadja el. Az ilyen gépet az $L$ nyelv \textbf{eldöntőjének} nevezzük.

\begin{summarybox}
Az eldöntésnél a megállás minden bemeneten kötelező. Nem elég az, hogy az igen esetekben előbb-utóbb választ kapunk; a nem esetekre is véges időben elutasítást kell adni.
\end{summarybox}

\subsection{Felismerhető probléma}

Egy $L$ nyelv \textbf{felismerhető}, ha létezik olyan Turing-gép, amely minden $w\in L$ bemenetre elfogadással megáll. Ha $w\notin L$, akkor a gép vagy elutasít, vagy soha nem áll meg.

Ez gyengébb feltétel, mint az eldönthetőség. A felismerő gép pozitív választ tud tanúsítani, de negatív választ nem feltétlenül ad.

\begin{center}
\begin{tabularx}{\textwidth}{L{31mm}L{45mm}X}
\toprule
\textbf{Fogalom} & \textbf{Mit várunk el a géptől?} & \textbf{Következmény} \\
\midrule
Elfogadás & Egy adott futás elfogadó állapotban megáll. & Egyetlen bemenetre vonatkozó viselkedés. \\
Felismerés & Ha $w\in L$, akkor a gép elfogadja $w$-t. Ha $w\notin L$, akkor elutasíthat vagy nem áll meg. & Az $L$ nyelv rekurzívan felsorolható. \\
Eldöntés & Minden $w$ bemeneten megáll, és elfogad pontosan akkor, ha $w\in L$. & Az $L$ nyelv rekurzív, azaz algoritmikusan eldönthető. \\
Felsorolás & Egy gép sorban kiírja az $L$ elemeit, ismétlés megengedett. & Nem üres nyelveknél ekvivalens a felismerhetőséggel. \\
\bottomrule
\end{tabularx}

\end{center}

\section{A Church-tézis}

\subsection{Az algoritmus intuitív és formális fogalma}

A hétköznapi algoritmusfogalom szerint algoritmusnak nevezünk egy végesen leírható, mechanikusan végrehajtható eljárást. Ez az intuitív meghatározás programozási gyakorlatban jól használható, de matematikai bizonyításokhoz nem elég pontos. Ha azt akarjuk bizonyítani, hogy valamire nincs algoritmus, akkor először pontosan meg kell mondani, mit értünk algoritmus alatt.

A 20. század első felében több, egymástól független formális számítási modell született:

\begin{itemize}
  \item Turing-gép;
  \item lambda-kalkulus;
  \item rekurzív függvények;
  \item később RAM-gép és más programozási modellek.
\end{itemize}

Ezekről a modellekről kiderült, hogy kiszámíthatósági értelemben ugyanazt tudják: ugyanazokat a függvényeket képesek kiszámítani, legfeljebb eltérő hatékonysággal.

\subsection{Church-tézis megfogalmazása}

A \textbf{Church-tézis} nem matematikai tétel, hanem elvi azonosítás:
\[
  \text{intuitívan algoritmikusan kiszámítható}
  \quad = \quad
  \text{Turing-géppel kiszámítható}.
\]

Más megfogalmazásban: minden olyan eljárás, amelyet ésszerűen algoritmusnak tekintünk, formalizálható Turing-géppel vagy azzal ekvivalens modellben.

Fontos, hogy ez nem bizonyítható hagyományos tételként, mert az egyik oldalon az „intuitív algoritmus” fogalma szerepel. A tézist az támasztja alá, hogy nagyon sok különböző, természetes számítási modell ugyanarra a kiszámítható függvényosztályra vezet.

\subsection{Jelentősége az eldönthetetlenségben}

A Church-tézis miatt, ha egy problémáról bebizonyítjuk, hogy Turing-géppel nem dönthető el, akkor ezt úgy értelmezzük, hogy nincs rá algoritmus semmilyen hagyományos, mechanikus értelemben.

Például a megállási probléma eldönthetetlensége nem csak azt jelenti, hogy „a Turing-gépek korlátozottak”. A jelentése sokkal erősebb: nincs általános program, amely minden programról és minden bemenetről eldöntené, hogy a program meg fog-e állni.

\section{Rekurzív és parciálisan rekurzív függvények}

\subsection{Teljes, kiszámítható függvény}

Legyen $f:\Sigma^*\to\Sigma^*$ egy függvény. Azt mondjuk, hogy $f$ \textbf{kiszámítható}, vagy klasszikus szóhasználattal \textbf{rekurzív}, ha van olyan Turing-gép, amely minden $x\in\Sigma^*$ bemeneten megáll, és kimenetként $f(x)$ értékét adja.

A lényeg a teljes megállás:
\[
  \forall x\in\Sigma^*:
  M(x) \text{ véges időben megáll, és } M(x)=f(x).
\]

Példák rekurzív függvényekre:

\begin{itemize}
  \item két természetes szám összeadása;
  \item egy szó megfordítása;
  \item egy szám prímtényezős felbontása, elméleti értelemben;
  \item egy véges automata szimulációja adott szón.
\end{itemize}

Ezeknél létezik olyan eljárás, amely minden megengedett bemeneten befejeződik.

\subsection{Parciálisan rekurzív függvény}

Egy $f:\Sigma^*\rightharpoonup\Sigma^*$ parciális függvény \textbf{parciálisan rekurzív}, ha van olyan Turing-gép, amely pontosan azokon a bemeneteken áll meg és ad kimenetet, ahol $f$ definiálva van. Ha $f(x)$ nincs definiálva, akkor a gép nem köteles megállni.

Formálisan:
\[
  f(x)\downarrow \Rightarrow M(x) \text{ megáll és } M(x)=f(x),
\]
\[
  f(x)\uparrow \Rightarrow M(x) \text{ nem áll meg, vagy nincs értelmezett kimenete.}
\]

A $\downarrow$ jel azt fejezi ki, hogy a függvényérték definiált, az $\uparrow$ pedig azt, hogy nem definiált.

\begin{summarybox}
A rekurzív függvény teljes: minden bemenetre ad választ. A parciálisan rekurzív függvény csak bizonyos bemeneteken köteles választ adni. Ez a különbség felel meg nyelveknél az eldöntés és felismerés különbségének.
\end{summarybox}

\section{Rekurzív és rekurzívan felsorolható nyelvek}

\subsection{Rekurzív nyelv}

Egy $L\subseteq\Sigma^*$ nyelv \textbf{rekurzív}, ha karakterisztikus függvénye rekurzív:
\[
  \chi_L(w)=
  \begin{cases}
    1, & \text{ha } w\in L,\\
    0, & \text{ha } w\notin L.
  \end{cases}
\]

Ez pontosan azt jelenti, hogy $L$ eldönthető. Van olyan Turing-gép, amely minden bemeneten megáll, és $1$ vagy $0$ választ ad.

Példák rekurzív nyelvekre:

\begin{itemize}
  \item páros hosszú szavak nyelve;
  \item reguláris nyelv szóproblémája;
  \item környezetfüggetlen nyelv szóproblémája;
  \item prímszámok megfelelő kódolású halmaza;
  \item bármely véges nyelv.
\end{itemize}

\subsection{Rekurzívan felsorolható nyelv}

Egy $L$ nyelv \textbf{rekurzívan felsorolható}, ha létezik olyan algoritmus, amely valamilyen sorrendben felsorolja az $L$ elemeit, vagy ekvivalens módon létezik olyan Turing-gép, amely pontosan az $L$-beli szavakon áll meg elfogadással.

A felismerő szemlélet:
\[
  w\in L \Rightarrow M(w) \text{ elfogadással megáll},
\]
\[
  w\notin L \Rightarrow M(w) \text{ elutasít vagy nem áll meg}.
\]

A felsorolási szemlélet szerint egy gép sorban kiírja az $L$ elemeit:
\[
  w_1,w_2,w_3,\ldots
\]
Ismétlés megengedett. Ha $L$ végtelen, akkor a felsorolás végtelen futás lehet; ez nem baj, mert a cél az, hogy minden $L$-beli elem előbb-utóbb megjelenjen.

\subsection{Felismerő pszeudókód}

A következő pszeudókód azt a tipikus felismerő működést mutatja, amely egy gépet szimulál. Ha a szimulált gép elfogad, a felismerő is elfogad. Ha a szimulált gép nem áll meg, akkor a felismerő sem jut eredményre.

\zvpseudocode[caption={Felismerés szimulációval},label={lst:zv11_felismero_szimulacio}]{src/code/zv11_felismero_szimulacio.pseudo}

\section{Az R, RE, coR és coRE osztályok}

\subsection{Az osztályok jelentése}

A tételben szereplő jelölések a következők:

\begin{itemize}
  \item $\mathrm{R}$: rekurzív, azaz eldönthető nyelvek osztálya;
  \item $\mathrm{RE}$: rekurzívan felsorolható, azaz felismerhető nyelvek osztálya;
  \item $\mathrm{coR}$: azok a nyelvek, amelyek komplementere $\mathrm{R}$-beli;
  \item $\mathrm{coRE}$: azok a nyelvek, amelyek komplementere $\mathrm{RE}$-beli.
\end{itemize}

Mivel a rekurzív nyelvek zártak komplementerre, ezért:
\[
  \mathrm{coR}=\mathrm{R}.
\]

Az $\mathrm{RE}$ osztály viszont általában nem zárt komplementerre. A megállási probléma tipikus példa:
\[
  \mathrm{HALT}\in\mathrm{RE},
  \qquad
  \mathrm{HALT}\notin\mathrm{R},
  \qquad
  \overline{\mathrm{HALT}}\notin\mathrm{RE}.
\]

\begin{figure}[htbp]
\centering
\begin{tikzpicture}[font=\small]
  \tikzset{mainlabel/.style={font=\bfseries, align=center},
           smallnote/.style={font=\scriptsize, align=center},
           rbox/.style={draw, rounded corners, fill=white, minimum width=31mm, minimum height=11mm, align=center},
           outbox/.style={draw, rounded corners, fill=black!4, minimum width=55mm, minimum height=9mm, align=center}}

  \draw[thick, fill=lightaccent, fill opacity=0.55] (-1.55,0) ellipse (2.45cm and 1.45cm);
  \draw[thick, fill=warnbg, fill opacity=0.55] (1.55,0) ellipse (2.45cm and 1.45cm);

  \node[mainlabel] at (-2.45,0.85) {$\mathrm{RE}$};
  \node[smallnote] at (-2.45,0.45) {felismerhető\\nyelvek};
  \node[mainlabel] at (2.45,0.85) {$\mathrm{coRE}$};
  \node[smallnote] at (2.45,0.45) {komplementere\\felismerhető};

  \node[rbox] at (0,0) {$\mathrm{R}=\mathrm{RE}\cap\mathrm{coRE}$\\eldönthető};

  \node[smallnote] at (-2.55,-0.75) {példa: $\mathrm{HALT}$\\$\mathrm{RE}$, de nem $\mathrm{R}$};
  \node[smallnote] at (2.55,-0.75) {példa: $\overline{\mathrm{HALT}}$\\$\mathrm{coRE}$, de nem $\mathrm{R}$};
  \node[outbox] at (0,-2.2) {vannak nyelvek, amelyek sem $\mathrm{RE}$-, sem $\mathrm{coRE}$-beliek};
\end{tikzpicture}
\caption{A legfontosabb kiszámíthatósági nyelvosztályok kapcsolata}
\label{fig:zv11_nyelvosztalyok}
\end{figure}


\subsection{Alapvető tartalmazások}

Minden eldönthető nyelv felismerhető:
\[
  \mathrm{R}\subseteq\mathrm{RE}.
\]
Ha ugyanis egy eldöntő gép minden bemeneten megáll, akkor felismerőként is használható: az igen válasz esetén elfogad, a nem válasz esetén elutasít.

A tartalmazás valódi:
\[
  \mathrm{R}\subsetneq\mathrm{RE}.
\]
Ennek klasszikus oka a megállási probléma: felismerhető, mert a gépet lehet szimulálni, és ha megáll, ezt észrevesszük; viszont nem dönthető el általánosan, mert nincs minden bemeneten megálló halting-teszt.

\subsection{A legfontosabb kapcsolat}

A rekurzív és rekurzívan felsorolható nyelvek között az egyik legfontosabb tétel:
\[
  L\in\mathrm{R}
  \quad\Longleftrightarrow\quad
  L\in\mathrm{RE}\text{ és }\overline L\in\mathrm{RE}.
\]
Másképpen:
\[
  \mathrm{R}=\mathrm{RE}\cap\mathrm{coRE}.
\]

Az egyik irány egyszerű: ha $L$ eldönthető, akkor $L$ is és $\overline L$ is felismerhető. A másik irány lényege, hogy ha van felismerőnk az igen esetekre és a nem esetekre is, akkor a két felismerőt váltogatva futtatjuk. Mivel minden bemenet vagy $L$-ben van, vagy a komplementerében, az egyik felismerő biztosan elfogad véges időben.

\begin{figure}[htbp]
\centering
\begin{tikzpicture}[font=\small]
  \tikzset{box/.style={draw, rounded corners, fill=lightaccent, minimum width=34mm, minimum height=10mm, align=center},
           good/.style={draw, rounded corners, fill=warnbg, minimum width=34mm, minimum height=10mm, align=center},
           arr/.style={->, >=stealth, thick},
           lab/.style={font=\scriptsize, align=center}}

  \node[box] (w) at (0,1.8) {bemenet $w$};
  \node[box] (ml) at (-2.5,0.35) {$M_L$ felismerő\\$L$-re};
  \node[box] (mc) at (2.5,0.35) {$M_{\overline L}$ felismerő\\$\overline L$-re};
  \node[good] (yes) at (-2.5,-1.25) {$M_L$ elfogad\\$\Rightarrow w\in L$};
  \node[good] (no) at (2.5,-1.25) {$M_{\overline L}$ elfogad\\$\Rightarrow w\notin L$};
  \node[lab] at (0,-2.25) {A két felismerőt váltogatva futtatva valamelyik biztosan megáll,\newline ezért $L$ eldönthető.};

  \draw[arr] (w) -- (ml);
  \draw[arr] (w) -- (mc);
  \draw[arr] (ml) -- (yes);
  \draw[arr] (mc) -- (no);
\end{tikzpicture}
\caption{Ha $L$ és komplementere is felismerhető, akkor $L$ eldönthető}
\label{fig:zv11_eldontes_es_felismeres}
\end{figure}


\zvpseudocode[caption={Eldöntő gép két felismerőből},label={lst:zv11_eldonto_ket_felismerobol}]{src/code/zv11_eldonto_ket_felismerobol.pseudo}

A pszeudókódban azért nem egyszerűen egymás után futtatjuk a két felismerőt, mert az első felismerő lehet, hogy nem áll meg. Ezért kell váltogatott vagy párhuzamos szimuláció: először egy lépést, majd kettőt, majd hármat engedünk mindkettőnek, és így tovább.

\subsection{Zártsági tulajdonságok}

A zártsági tulajdonságok vizsgán azért hasznosak, mert gyorsan megmutatják, milyen műveletekkel maradunk egy adott nyelvosztályon belül.

\begin{center}
\begin{tabularx}{\textwidth}{L{34mm}L{37mm}X}
\toprule
\textbf{Tulajdonság} & \textbf{Rekurzív nyelvek} & \textbf{Rekurzívan felsorolható nyelvek} \\
\midrule
Unió & Zártak: párhuzamosan vagy egymás után dönthetők. & Zártak: ha bármelyik felismerő elfogad, elfogadunk. \\
Metszet & Zártak: mindkét eldöntőt lefuttatjuk. & Zártak: mindkét felismerő elfogadását ki lehet várni. \\
Komplementer & Zártak: az igen/nem választ felcseréljük. & Nem zártak általában. Ha $L$ és $\overline L$ is $\mathrm{RE}$, akkor $L\in\mathrm{R}$. \\
Különbség & Zártak: $L_1\setminus L_2=L_1\cap\overline{L_2}$. & Nem zártak általában, mert komplementerre lenne szükség. \\
\bottomrule
\end{tabularx}

\end{center}

A legfontosabb különbség a komplementer. Rekurzív nyelvnél az eldöntő gép igen/nem válaszát fel lehet cserélni. Rekurzívan felsorolható nyelvnél viszont lehet, hogy a felismerő a nem esetekben soha nem áll meg, ezért nincs mit egyszerűen „megfordítani”.

\section{A megállási probléma és eldönthetetlenség}

\subsection{A megállási probléma}

A megállási probléma nyelve:
\[
  \mathrm{HALT}=\{\langle M,w\rangle\mid M \text{ megáll } w \text{ bemeneten}\}.
\]

Ez a nyelv rekurzívan felsorolható. A felismerő algoritmus egyszerű: szimuláljuk $M$ futását $w$-n. Ha $M$ megáll, akkor a szimuláció ezt véges időben észreveszi, és elfogad.

A döntéshez viszont azt is meg kellene mondani véges időben, ha $M$ soha nem fog megállni. Általánosan ez lehetetlen.

\subsection{Diagonális bizonyítás gondolata}

Tegyük fel indirekt módon, hogy létezik egy $H$ gép, amely minden $\langle M,w\rangle$ párra eldönti, hogy $M$ megáll-e $w$-n. Ebből készítünk egy $D$ gépet, amely egy gépleírást kap bemenetként, és önmagára alkalmazza a vizsgált gépet:

\begin{itemize}
  \item ha $H$ szerint $M$ megáll $\langle M\rangle$ bemeneten, akkor $D$ direkt végtelen ciklusba megy;
  \item ha $H$ szerint $M$ nem áll meg $\langle M\rangle$ bemeneten, akkor $D$ megáll.
\end{itemize}

Ekkor megkérdezhetjük: mi történik $D(\langle D\rangle)$ futásakor? Ha $H$ szerint megáll, akkor $D$ úgy van definiálva, hogy nem áll meg. Ha $H$ szerint nem áll meg, akkor $D$ megáll. Mindkét eset ellentmondás, tehát a feltételezett $H$ döntő nem létezhet.

\begin{figure}[htbp]
\centering
\begin{tikzpicture}[font=\small]
  \tikzset{box/.style={draw, rounded corners, fill=lightaccent, minimum width=34mm, minimum height=9mm, align=center},
           warn/.style={draw, rounded corners, fill=warnbg, minimum width=36mm, minimum height=9mm, align=center},
           arr/.style={->, >=stealth, thick},
           lab/.style={font=\scriptsize, align=center}}

  \node[box] (in) at (-3.9,0) {bemenet\\$\langle M\rangle$};
  \node[warn] (d) at (0,0) {diagonális\\gép $D$};
  \node[box] (h) at (0,1.55) {feltételezett\\döntő $H$};
  \node[box] (query) at (3.9,0) {$H(\langle M\rangle,\langle M\rangle)$\\kérdés};
  \node[warn] (contr) at (0,-1.55) {$D(\langle D\rangle)$\\ellentmondás};

  \draw[arr] (in) -- (d);
  \draw[arr] (d) -- (query);
  \draw[arr] (query) -- (h);
  \draw[arr] (h) -- (d);
  \draw[arr] (d) -- (contr);

  \node[lab] at (3.9,-0.82) {önmagára alkalmazás};
\end{tikzpicture}
\caption{A megállási probléma eldönthetetlenségének diagonális gondolata}
\label{fig:zv11_megallasi_diagonalizalas}
\end{figure}


\zvpseudocode[caption={A diagonális gép gondolata},label={lst:zv11_diagonalis_gep}]{src/code/zv11_diagonalis_gep.pseudo}

\begin{summarybox}
A diagonális bizonyítás lényege az önalkalmazás. A feltételezett döntőgépből olyan gépet építünk, amely saját magára alkalmazva pontosan az ellenkezőjét csinálja annak, amit a döntő jósol. Ez ellentmondást ad.
\end{summarybox}

\section{Nevezetes eldönthetetlen problémák}

\subsection{Megállási és elfogadási problémák}

A megállási probléma mellett gyakran szerepel az elfogadási probléma:
\[
  \mathrm{A}_{TM}=\{\langle M,w\rangle\mid M \text{ elfogadja } w\text{-t}\}.
\]
Ez is felismerhető: szimuláljuk $M$-et $w$-n, és ha elfogad, mi is elfogadunk. Ugyanakkor nem eldönthető, mert ha eldönthető lenne, akkor a megállási probléma is hasonló kódolással eldönthetővé válna.

A nem-megállás nyelve:
\[
  \overline{\mathrm{HALT}}=
  \{\langle M,w\rangle\mid M \text{ nem áll meg } w\text{-n}\}
\]
nem rekurzívan felsorolható. Ha az lenne, akkor $\mathrm{HALT}$ és komplementere is $\mathrm{RE}$-beli lenne, így $\mathrm{HALT}$ rekurzív lenne, ami ellentmondás.

\begin{center}
\small
\begin{tabularx}{\textwidth}{L{32mm}L{38mm}X}
\toprule
\textbf{Probléma} & \textbf{Nyelv alakja} & \textbf{Besorolás} \\
\midrule
Megállási probléma & $\mathrm{HALT}$: $M$ megáll $w$-n & $\mathrm{RE}$, de nem $\mathrm{R}$. Szimulációval felismerhető, de általánosan nem eldönthető. \\
Elfogadási probléma & $\mathrm{A}_{TM}$: $M$ elfogadja $w$-t & $\mathrm{RE}$, de nem $\mathrm{R}$. Ha a szimuláció elfogadást lát, megállhat. \\
Nem-megállás & $\overline{\mathrm{HALT}}$ & $\mathrm{coRE}$, de nem $\mathrm{RE}$. Ha ez is $\mathrm{RE}$ lenne, akkor $\mathrm{HALT}$ eldönthető lenne. \\
Üres-e a felismert nyelv? & $L(M)=\emptyset$ & Eldönthetetlen; tipikusan nem is $\mathrm{RE}$. Rice-tétel jellegű következmény. \\
Reguláris nyelv szóproblémája & Véges automata elfogadja-e $w$-t & $\mathrm{R}$, sőt hatékonyan eldönthető. A véges automata futása mindig véges. \\
\bottomrule
\end{tabularx}
\normalsize

\end{center}

\subsection{Rice-tétel szemlélete}

A Rice-tétel azt mondja ki, hogy a Turing-gépek által felismert nyelvek minden nemtriviális szemantikus tulajdonsága eldönthetetlen.

Szemantikus tulajdonság például:

\begin{itemize}
  \item üres-e a gép által felismert nyelv;
  \item véges-e a felismert nyelv;
  \item reguláris-e a felismert nyelv;
  \item tartalmaz-e egy adott szót;
  \item minden szót elfogad-e a gép.
\end{itemize}

Nemtriviális azt jelenti, hogy a tulajdonság nem igaz minden Turing-felismerhető nyelvre, és nem is hamis mindegyikre. A tétel azért erős, mert nem csak egy-egy konkrét problémát tilt, hanem egész problémacsaládokról mondja ki, hogy eldönthetetlenek.

\begin{warningbox}
A Rice-tétel a felismert nyelv tulajdonságaira vonatkozik, nem a gép puszta szintaktikai jellemzőire. Például az, hogy egy gép leírásában hány állapot szerepel, szintaktikai kérdés és eldönthető. Az, hogy a gép pontosan milyen nyelvet ismer fel, szemantikai kérdés.
\end{warningbox}

\section{Redukciók szerepe eldönthetetlenségi bizonyításban}

\subsection{Az alapötlet}

Eldönthetetlenséget gyakran redukcióval bizonyítunk. Ha ismert, hogy egy $A$ probléma eldönthetetlen, és meg tudjuk mutatni, hogy $A$ megoldható lenne egy $B$ probléma döntőjének segítségével, akkor $B$ sem lehet eldönthető.

Szimbolikusan:
\[
  A \leq_m B,
\]
ahol ez azt jelenti, hogy az $A$ problémából hatékonyan vagy legalább algoritmikusan át tudunk alakítani példányokat a $B$ problémába úgy, hogy az igen/nem válasz megmaradjon.

Ha $A$ eldönthetetlen és $A\leq_m B$, akkor $B$ eldönthetetlen. Ellenkező esetben ugyanis $B$ döntője és az átalakítás együtt $A$ döntőjét adná.

\subsection{Tipikus bizonyítási séma}

Egy eldönthetetlenségi bizonyítás gyakori felépítése:

\begin{enumerate}
  \item kiválasztunk egy már ismert eldönthetetlen problémát, például $\mathrm{HALT}$-ot;
  \item feltesszük, hogy a vizsgált $B$ probléma eldönthető;
  \item megmutatjuk, hogyan lehetne ezzel $\mathrm{HALT}$-ot eldönteni;
  \item ez ellentmond a megállási probléma eldönthetetlenségének;
  \item ezért $B$ is eldönthetetlen.
\end{enumerate}

Ez a gondolat jelenik meg például annak bizonyításában, hogy nem dönthető el általánosan, üres-e egy Turing-gép által felismert nyelv, vagy elfogad-e a gép legalább egy szót.

\section{Polinomiális idejű algoritmusok}

\subsection{Miért kell bonyolultságot vizsgálni?}

Az eldönthetőség csak azt kérdezi, hogy létezik-e véges időben befejeződő algoritmus. Ez azonban gyakorlati szempontból kevés. Egy algoritmus lehet eldöntő, de olyan lassú, hogy nagyobb bemeneteken használhatatlan.

Ezért vezetjük be az időbonyolultságot. Egy Turing-gép időigénye:
\[
  T_M(n)=\max\{\text{$M$ lépésszáma $w$-n}\mid |w|=n\}.
\]
Ez a legrosszabb esetű futási idő a bemenet hossza szerint.

\subsection{Polinomiális idő}

Egy algoritmus \textbf{polinomiális idejű}, ha létezik olyan konstans $c$, amelyre:
\[
  T_M(n)=O(n^c).
\]

A polinomiális időben eldönthető nyelvek osztálya:
\[
  \mathrm{P}=\bigcup_{c\geq 1}\mathrm{DTIME}(n^c).
\]

A $\mathrm{P}$ osztály az eldönthető problémák hatékonyan kezelhető részét jelöli. Minden $\mathrm{P}$-beli nyelv rekurzív:
\[
  \mathrm{P}\subseteq\mathrm{R},
\]
hiszen ha egy gép polinomiális időben eldönt egy nyelvet, akkor minden bemeneten véges időben megáll.

\begin{center}
\begin{tabularx}{\textwidth}{L{31mm}L{44mm}X}
\toprule
\textbf{Fogalom} & \textbf{Definíció / jelölés} & \textbf{Értelmezés} \\
\midrule
Időigény & $T_M(n)$ a legrosszabb futási idő $n$ hosszú bemeneteken. & A bemenethossz függvényében mérjük, nem egy konkrét bemenetre. \\
Polinomiális idő & $T_M(n)=O(n^c)$ valamely konstans $c$ mellett. & Elméleti értelemben hatékonynak tekintett algoritmus. \\
$\mathrm{P}$ osztály & Polinomiális időben eldönthető nyelvek osztálya. & $\mathrm{P}\subseteq\mathrm{R}$, mert polinomiális időben futó eldöntő gép biztosan megáll. \\
Exponenciális idő & Például $O(2^n)$ vagy $O(c^n)$. & Sok kombinatorikus keresési probléma naiv megoldása ilyen. \\
Kódolás szerepe & A bemenet hossza a leírás hosszát jelenti. & Egy szám értéke és a számjegyeinek száma nagyon eltérő lehet, ezért fontos a bitméret szerinti elemzés. \\
\bottomrule
\end{tabularx}

\end{center}

\subsection{Példák polinomiális algoritmusokra}

Polinomiális időben eldönthető példák:

\begin{itemize}
  \item egy szó illeszkedik-e egy reguláris kifejezéshez vagy véges automatához;
  \item van-e út két csúcs között egy gráfban;
  \item összefüggő-e egy gráf;
  \item minimális feszítőfa meghatározása;
  \item legrövidebb utak bizonyos feltételek mellett;
  \item prímszám-e egy adott binárisan kódolt természetes szám.
\end{itemize}

Nem minden eldönthető probléma ismert polinomiálisnak. Vannak eldönthető, de nagyon nagy időigényű problémák, és vannak olyanok is, amelyekre a legjobb ismert algoritmusok exponenciálisak.

\subsection{A bemenethossz szerepe}

Bonyolultságelméletben a bemenet mérete a kódolás hossza. Ez különösen számoknál fontos. Ha egy természetes szám értéke $N$, akkor bináris kódolásban a hossza körülbelül $\log_2 N$. Emiatt nem ugyanaz azt mondani, hogy egy algoritmus polinomiális $N$ értékében, vagy polinomiális a bemenet bitméretében.

Például egy olyan algoritmus, amely $O(N)$ lépést végez egy $N$ értékű számra, a bitméret szerint exponenciális is lehet, mert $N\approx 2^n$, ahol $n$ a bináris hossz.

\zvpseudocode[caption={Polinomiális időkorlát szemlélete},label={lst:zv11_polinomialis_eldonto}]{src/code/zv11_polinomialis_eldonto.pseudo}

\section{Kapcsolat az eldönthetőség és a bonyolultság között}

Az eldönthetőség és a bonyolultság két külön szint:

\begin{itemize}
  \item eldönthetőség: van-e egyáltalán minden bemeneten megálló algoritmus;
  \item bonyolultság: ha van, akkor mennyi időt vagy tárat igényel.
\end{itemize}

A kapcsolat szemléletesen:
\[
  \mathrm{P}\subseteq\mathrm{R}\subsetneq\mathrm{RE}\subseteq\mathcal P(\Sigma^*).
\]

Itt $\mathcal P(\Sigma^*)$ az összes nyelv halmaza. Nem minden nyelv $\mathrm{RE}$-beli, mert csak megszámlálhatóan sok Turing-gép van, de $\Sigma^*$ részhalmazaiból nem megszámlálhatóan sok van.

A vizsgán fontos hangsúlyozni, hogy az eldönthetetlen problémák nem azért rosszak, mert lassú algoritmusunk van rájuk, hanem mert nincs rájuk minden bemeneten működő algoritmus. A nehéz, de eldönthető problémáknál viszont van algoritmus, csak annak erőforrásigénye lehet túl nagy.

\section{Tipikus vizsgahibák és pontosítások}

\subsection{Felismerhető nem ugyanaz, mint eldönthető}

Gyakori hiba azt mondani, hogy ha egy gép felismer egy nyelvet, akkor el is dönti. Ez csak akkor igaz, ha a gép minden bemeneten megáll. A felismerő gép a nem elemekre végtelen futásba is kerülhet.

Helyes megfogalmazás:
\[
  \text{eldönthető} \Rightarrow \text{felismerhető},
\]
de visszafelé általában nem igaz.

\subsection{A komplementer problémája}

Ha $L\in\mathrm{RE}$, abból nem következik, hogy $\overline L\in\mathrm{RE}$. Ez a különbség az egyik legfontosabb oka annak, hogy $\mathrm{RE}$ nagyobb, mint $\mathrm{R}$.

Ha viszont mindkettő igaz:
\[
  L\in\mathrm{RE}\quad\text{és}\quad\overline L\in\mathrm{RE},
\]
akkor:
\[
  L\in\mathrm{R}.
\]

\subsection{Church-tézis nem formális tétel}

A Church-tézis nem olyan állítás, amelyet a Turing-gép definíciójából bizonyítani lehetne. Azért tézis, mert az algoritmus intuitív fogalmát azonosítja a formális kiszámíthatósággal. Vizsgán érdemes így megfogalmazni:

\begin{quote}
A Church-tézis szerint minden intuitív értelemben algoritmikusan kiszámítható függvény Turing-géppel is kiszámítható.
\end{quote}

\subsection{Polinomiális idő nem azonos az eldönthetőséggel}

A polinomiális idő már bonyolultsági kategória. Minden polinomiális időben eldönthető probléma eldönthető, de nem minden eldönthető probléma polinomiális idejű. Az eldönthetetlen problémák pedig nem tartoznak ide, mert ott nincs minden bemeneten megálló döntő algoritmus.

\section{Összefoglalás}

\begin{summarybox}
A 11. tétel központi gondolata, hogy a pontosan megfogalmazott problémák között is vannak olyanok, amelyekre nem létezik általános döntő algoritmus. Az $\mathrm{R}$ osztály az eldönthető nyelveket, az $\mathrm{RE}$ osztály a felismerhető nyelveket tartalmazza. A rekurzív nyelvek zártak komplementerre, ezért $\mathrm{coR}=\mathrm{R}$, de az $\mathrm{RE}$ osztály általában nem zárt komplementerre. A megállási probléma $\mathrm{RE}$-beli, de nem rekurzív, ezért alapvető példa eldönthetetlen problémára. A polinomiális idejű algoritmusok az eldönthető problémák hatékonyan megoldható részét adják.
\end{summarybox}

A tétel végén érdemes fejben tartani a következő láncot:
\[
  \mathrm{P}\subseteq\mathrm{R}
  =\mathrm{RE}\cap\mathrm{coRE}
  \subsetneq\mathrm{RE}.
\]

Rövid vizsgaválaszban a megállási problémát mindenképpen érdemes megemlíteni, mert ezen keresztül jól látszik a felismerés és az eldöntés különbsége: ha a program megáll, azt szimulációval észrevesszük, de ha nem áll meg, nincs általános módszerünk ennek véges időben történő felismerésére.
